\chapter{紧线性映射与扰动}
\setcounter{exercisepar}{0}

在本章中，所考虑的向量空间都是等于 R 或 C 的域 K 上的向量空间。本书中所作的关于 EVT, II 的引用一般针对 K = R 的情形；关于 K = C 的情形，参见 EVT, II, p. 64 至 68。

半赋范空间是指赋有一个半范数 p 以及由 p 定义的拓扑的向量空间 E（参见 EVT, II, p. 2）。E 的单位球，或 p 的单位球，是 E 中满足 \(p(x) \leqslant 1\) 的元素 x 所成的集。

对每个向量空间 E，我们用 \(1_{E}\) 表示 E 的恒等映射。若 E, F 和 G 是向量空间，且 u: \(E \rightarrow F\)、\(v: F \rightarrow G\) 是线性映射，我们有时用 vu 表示从 E 到 G 的线性映射 \(v \circ u\)。

设 E 和 F 是拓扑向量空间。从 E 到 F 的连续线性映射也称为态射，而当 F = E 时称为自同态。从 E 到 F 的双射线性映射，若它连续且其逆也连续，则称为同构，而当 F = E 时称为自同构。当 E 和 F 是半赋范空间时，自同态、同构和自同构这几个概念都是就底层的拓扑向量空间结构而言的。

给定拓扑向量空间 E 和 F，我们用 \(\mathcal{L}^{\mathrm{f}}(\mathrm{E};\mathrm{F})\) 表示从 E 到 F 的有限秩连续线性映射所成的向量空间。我们也用 \(\mathcal{L}^{\mathrm{f}}(\mathrm{E})\) 表示向量空间 \(\mathcal{L}^{\mathrm{f}}(\mathrm{E};\mathrm{E})\)。

我们回顾，从 E 到 F 的连续线性映射 u 称为严格的，如果它通过过渡到商空间诱导出从 E/ Ker(u) 到 u(E) 上的一个同构（TG, III, p. 16, def. 1）；这等价于说，E 中 0 的每个邻域的像是 u(E) 中 0 的一个邻域（TG, III, p. 16, prop. 24）。

